\documentclass[11pt]{article}
\usepackage[a4paper,margin=1in]{geometry}
\usepackage{amsmath,amssymb}
\usepackage{enumitem}
\usepackage[T1]{fontenc}
\usepackage[utf8]{inputenc}

\title{SRT Valuation Model — Locked Parameters (Sections 1--8)}
\date{2026-02-28}

\begin{document}
\maketitle

\begin{enumerate}[label=\arabic*),leftmargin=*,align=left]
\item \textbf{As-of anchor}: Valuation is anchored to the market/curve date.

\item \textbf{As-of timestamp convention}: As-of is end-of-day (COB) on the as-of date.

\item \textbf{“Future” cashflows inclusion rule}: Only include cashflows with cashflow date > as-of date (strictly greater; if equal to as-of date it is excluded).

\item \textbf{Premium payment frequency}: Quarterly.

\item \textbf{Payment date sourcing}: Rules-only (no hybrid list+rules); schedule is generated from rules + user inputs.

\item \textbf{Schedule anchor input}: User must input First Payment Date (calendar date).

\item \textbf{First Payment Date validation}: If the user enters an invalid calendar date -> hard error.

\item \textbf{Quarterly stepping rule}: Generate dates by adding +3 calendar months each step (not 0.25y).

\item \textbf{EOM toggle required}: User must explicitly choose EOM = ON/OFF (no default).

\item \textbf{EOM behavior (if ON)}: If the anchor date is month-end, generated dates stay month-end (e.g., Jan 31 -> Apr 30 -> Jul 31).

\item \textbf{EOM behavior (if OFF)}: Keep same day-of-month where possible; if the day doesn’t exist in a month, use the last valid day of that month (but don’t force month-end thereafter).

\item \textbf{Business day adjustment for payment dates}: Use Modified Following.

\item \textbf{Calendar selection method}: User chooses the holiday calendar from a dropdown.

\item \textbf{Allowed calendars (dropdown set)}: TARGET/TARGET2; Sweden; UK; Germany; France; Spain; Switzerland.

\item \textbf{Joint calendar support}: User can choose Joint of N calendars; a date is a business day only if it is a business day in all selected calendars.

\item \textbf{Business day definition}: Non-business days include weekends AND holidays per the selected calendar(s).

\item \textbf{Effective accrual start for valuation}: We do not assume/compute past accrual. Accrual in valuation starts at $t_{\text{start,eff}} = \max(\text{as-of date}, \text{Accrual Start Date})$.

\item \textbf{Accrual Start Date input}: User must input Accrual Start Date (required, because it may be after as-of).

\item \textbf{Accrual End Date input}: User inputs Accrual End Date; if unknown, set Accrual End = Protection End (see later items).

\item \textbf{Accrual date validation}: If Accrual Start Date > Accrual End Date -> hard error.

\item \textbf{First payment date selection anchor}: The first generated payment date is the first payment date strictly after $t_{\text{start,eff}}$ (not strictly after as-of).

\item \textbf{Future-only schedule generation (valuation)}: Generate payment dates forward only, starting from the first payment date strictly after $t_{\text{start,eff}}$, then +3M until Accrual End.

\item \textbf{Final stub payment date rule}: If Accrual End is not already a scheduled payment date, add Accrual End as a final stub payment date (then apply business-day adjustment rules to that date).

\item \textbf{Final stub vs protection end}: The payment date may fall after Protection End; what matters is accrual stops at Accrual End / Protection End per rules (payment can occur after).

\item \textbf{Accrual interval inclusion convention}: Use include start, exclude end for year-fractions and accrual windows: [start, end).

\vspace{1em}
\noindent Locked parameters 26–50
\vspace{0.5em}

\item \textbf{Premium day count input method}: User selects the day count convention for premium accrual from a dropdown.

\item \textbf{Premium day count default}: If user doesn’t know, default day count = 30E/360.

\item \textbf{Day count dropdown set}: 30E/360; ACT/360; ACT/365F; 30/360 (US); ACT/ACT (ISDA); ACT/ACT (ICMA).

\item \textbf{Premium amount formula (generic)}: For any accrual window $[a,b)$, premium cash due is computed actuarially as:
\[ Premium(a,b) = s \cdot \text{YearFrac}(a,b) \cdot (\text{premium base over } [a,b)) \]
where $s$ is the premium spread (input elsewhere) and the premium base may change at event dates (see notice step-down rules).

\item \textbf{Premium cash timing}: Premium is paid on the payment date (no settlement lag).

\item \textbf{Accrual end governs accrual stop}: Premium accrues until Effective Accrual End Date; premium can still be paid after accrual ends (on the next payment date / final stub payment date).

\item \textbf{Protection End Date input}: User must input Protection End Date (required).

\item \textbf{Legal Final Maturity Date input}: User must input Legal Final Maturity Date (required).

\item \textbf{Protection End date adjustment}: Protection End is business-day adjusted using Modified Following + selected calendar(s).

\item \textbf{Legal Final date adjustment}: Legal Final is business-day adjusted using Modified Following + selected calendar(s).

\item \textbf{Both key dates required}: Protection End and Legal Final are both mandatory; missing/invalid -> hard error.

\item \textbf{Legal Final usage}: Legal Final is a boundary/cap only (no automatic “final cashflow event” node is created just because it’s Legal Final).

\item \textbf{Protection Start Date input}: User must input Protection Start Date (required).

\item \textbf{Protection Start date adjustment}: Protection Start is business-day adjusted using Modified Following + selected calendar(s).

\item \textbf{Coverage eligibility test (default-date based)}: A default is covered iff $\text{ProtStart} \le \tau \le \text{ProtEnd}$, with inclusive boundaries and using Default Date $\tau$ (not notice date) for eligibility.

\item \textbf{Notice may occur after Protection End}: If $\tau$ is covered, the write-down can occur at Notice Date even if Notice Date > ProtEnd (still capped by Legal Final).

\item \textbf{Notice lag definition}: $\text{NoticeDateRaw} = \tau + 1 \text{ calendar month}$ (and this is fixed; not user-configurable).

\item \textbf{Notice +1M month-add edge rule}: If the resulting day doesn’t exist (e.g., Jan 31 + 1M), clamp to month-end (Feb 28/29).

\item \textbf{Notice business-day adjustment}: If NoticeDateRaw is not a business day, adjust using Preceding (move to the previous business day) with the selected calendar(s) (including joint-of-N).

\item \textbf{Adjust-then-cap ordering}: Adjust Legal Final using Modified Following. Adjust NoticeDateRaw using Preceding. Then cap: $\text{NoticeDate} = \min(\text{AdjustedNotice}, \text{AdjustedLegalFinal})$.

\item \textbf{Legal Final cap rule}: Notice Date is never later than Legal Final after the above cap.

\item \textbf{Write-down timing intraday}: Write-down is effective at start-of-day on the (adjusted/capped) Notice Date.

\item \textbf{Premium accrual vs Notice Date}: Premium accrues up to but excluding Notice Date for affected exposures; the step-down happens at start-of-day Notice Date.

\item \textbf{Premium stop date for defaulted exposure}: Premium for the defaulted exposure stops at Notice Date (not at Default Date $\tau$).

\item \textbf{Premium base during Default->Notice lag}: From Default Date until Notice Date, premium is charged on the full pre-write-down notional (the notional is only reduced when written down at Notice).

\vspace{1em}
\noindent Locked parameters 51–75
\vspace{0.5em}

\item \textbf{Write-down amount formula (loan-level)}: On Notice Date, write-down for a defaulted loan is:
\[ WD_{\text{loan}} = EAD_{\text{loan}}(\tau) \cdot \max(LGD_{\text{recon}}, LGD_{\text{reg}}) \]

\item \textbf{LGD sources}: lgd\_recon and lgd\_reg are loan-level columns in the tape.

\item \textbf{LGD validation}: Hard validate for every loan: $0 \le LGD_{\text{recon}} \le 1$ and $0 \le LGD_{\text{reg}} \le 1$. Missing/invalid -> hard error.

\item \textbf{EAD timing for write-down}: EAD is evaluated at Default Date $\tau$ and frozen (so EAD at Notice = EAD at Default).

\item \textbf{EAD freeze at default}: Once an obligor defaults, its loans’ balances freeze at Default Date (no further amortization after $\tau$; notice lag doesn’t change EAD).

\item \textbf{Premium within-period handling}: Premium is computed by exact step-down integration at each Notice Date (piecewise-constant notional segments; premium base changes at Notice start-of-day).

\item \textbf{Notice Date adjustment consistency}: The adjusted Notice Date (after Preceding and capping) is used consistently for both write-down timing and premium-stop/step-down timing.

\item \textbf{Defaults after Protection End not covered}: If $\tau > \text{Protection End}$, the default is not covered (no notice/write-down triggered by it).

\item \textbf{Protection End equality}: If $\tau = \text{Protection End}$, it is covered (inclusive boundary).

\item \textbf{Loan maturity beyond Legal Final}: Allowed (loan maturity dates may extend past deal Legal Final; deal cashflows are still capped by Legal Final).

\item \textbf{Loan maturity date validity}: Missing/invalid maturity date -> hard error (unless explicitly excluded by rule “$T \le \text{as-of}$ is ignored”, see later).

\item \textbf{Tape currency assumption}: Outstanding\_Principal\_Amount is assumed to already be in deal currency. If not, hard error (no FX conversion in v1).

\item \textbf{Starting notional source column}: The tape contains Outstanding\_Principal\_Amount as the as-of outstanding balance per loan.

\item \textbf{Reference notional at as-of}: 
\[ N_{\text{ref}} = \sum_{\text{loans}} \text{Outstanding\_Principal\_Amount} \]
(no override/reconciliation input; tape sum is truth).

\item \textbf{Tranche sizing input}: User inputs tranche\_percentage (e.g., 6\%), meaning tranche size as a \% of $N_{\text{ref}}$.

\item \textbf{Our holding input}: User inputs our\_percentage, meaning our participation/holding share of the tranche.

\item \textbf{Our starting notional / premium base at as-of}: 
\[ N_0 = N_{\text{ref}} \cdot \text{tranche\_percentage} \cdot \text{our\_percentage} \]

\item \textbf{Loan-level write-down scaling to our position (v1)}: 
\[ WD_{\text{ours}} = \sum_{\text{defaulted\_loans}} WD_{\text{loan}} \cdot \text{tranche\_percentage} \cdot \text{our\_percentage} \]
(before attachment/detachment waterfall).

\item \textbf{Notional floor / cap}: Cumulative write-downs cannot exceed $N_0$. Remaining tranche notional: $N(t) = \max(N_0 - CumWD(t), 0)$.

\item \textbf{Position termination at zero}: If $N(t)=0$, the position is dead; it cannot go negative.

\item \textbf{Premium stops when notional hits zero}: If a Notice Date write-down causes $N(t)$ to hit 0, premium accrual stops immediately from that Notice Date start-of-day onward.

\item \textbf{No “accrued from past”}: We do not compute or assume any accrual before as-of. First accrual window in valuation begins at $t_{\text{start,eff}} = \max(\text{as-of}, \text{Accrual Start})$.

\item \textbf{First coupon proration scope}: Only the first coupon after valuation is prorated from $t_{\text{start,eff}}$ to the next payment date; later coupons accrue payment-to-payment.

\item \textbf{Payment dates strictly future}: Only payment dates with date > as-of are included for PV.

\item \textbf{Legal Final cap on notice/write-down}: If Notice Date (after +1M and adjustments) would exceed Legal Final, it is capped to Legal Final (after adjust-then-cap rule).

\vspace{1em}
\noindent Locked parameters 76–100
\vspace{0.5em}

\item \textbf{EAD projection method}: Rule-based using amortization type + maturity date from the tape.

\item \textbf{Amortization type mapping (tape -> model)}: Bullet -> bullet; n/a -> bullet; until further notice -> bullet; Amortisation -> linear rule; Annuity -> same linear rule.

\item \textbf{Linear balance rule for Amortisation/Annuity}: For as-of balance $B_0$, as-of date $t_0$, maturity $T$: 
\[ B(t) = B_0 \cdot \max(0, 1 - (t - t_0)/(T - t_0)) \]

\item \textbf{Linear rule time basis}: Use calendar-day ratio: $(t - t_0)/(T - t_0) = (\text{days}(t)-\text{days}(t_0))/(\text{days}(T)-\text{days}(t_0))$.

\item \textbf{Bullet balance rule}: $B(t) = B_0$ for $t < T$, and $B(t)=0$ for $t \ge T$.

\item \textbf{Loan maturity date adjustment}: Loan maturity $T$ is unadjusted calendar date (no business-day adjustment).

\item \textbf{Maturity timing intraday}: Balance drops to 0 at start-of-day on maturity date $T$.

\item \textbf{Exclude matured loans at/before as-of}: If a loan has $T \le \text{as-of date}$, ignore it (exclude from forward valuation).

\item \textbf{Exclude zero-balance loans}: If Outstanding\_Principal\_Amount = 0 at as-of, exclude the loan from forward valuation.

\item \textbf{Default vs maturity}: If simulated obligor default time $\tau > T$ for a given loan, treat that loan as effectively no default (out of risk set).

\item \textbf{Maturity reduces premium base}: Loan amortization/maturity reduces the portfolio balance and therefore reduces tranche notional/premium base over time (even without defaults).

\item \textbf{Amortization included in premium base}: Premium base reflects projected balances $B(t)$ over time and steps down at Notice Dates.

\item \textbf{Event-date evaluation}: Premium base/EAD evolution is evaluated on event dates, not daily.

\item \textbf{Event-date set includes}: at minimum include as-of date; payment dates; notice dates; accrual end date; loan maturity dates.

\item \textbf{Default time representation}: Simulated default time $\tau$ is treated as a calendar date (date-only; no time-of-day).

\item \textbf{tau -> date mapping}: Use the survival-curve implementation’s native time basis when converting/inverting to dates (“curve-native”).

\item \textbf{Curve horizon assumption}: Survival/hazard curves are guaranteed to extend beyond the SRT horizon (Protection End and Legal Final), so no extrapolation rule is needed.

\item \textbf{Marginal assignment by rating (loan-level)}: Each loan is assigned a rating bucket (via mapping) used in obligor curve selection.

\item \textbf{Obligor grouping key}: Group loans by Debtor\_ID (obligor-level modelling).

\item \textbf{Within-obligor linkage}: If obligor defaults, all loans under that Debtor\_ID default together (shared $\tau$).

\item \textbf{Default is absorbing}: An obligor can default at most once; no cure/re-default.

\item \textbf{Worst-rating obligor curve}: Each obligor uses the survival curve corresponding to the worst rating among its loans (replaces earlier weighted approach).

\item \textbf{Worst-rating ordering}: Use a fixed standard rating ordering (AAA best -> ... -> D worst).

\item \textbf{Internal rating source column}: Loan tape contains Internal\_Rating used for mapping to external rating bucket.

\item \textbf{Internal rating mapping}: Internal\_Rating is translated to a rating bucket using a provided mapping table (left -> right column).

\vspace{1em}
\noindent Locked parameters 101–125
\vspace{0.5em}

\item \textbf{Internal\_Rating exact-match rule}: Internal\_Rating must match a key in the mapping table exactly; otherwise hard error.

\item \textbf{Rating label normalization}: After mapping, normalize rating labels for lookup (trim whitespace, standardize case/format like “AA-”).

\item \textbf{Rating coverage requirement}: Mapped rating bucket must exist in the CDS curve set; otherwise hard error.

\item \textbf{Global data-quality stance (general)}: Missing/invalid/erroneous required data -> hard error (no silent fallbacks).

\item \textbf{Debtor\_ID missing}: Missing Debtor\_ID -> hard error (covered by global rule).

\item \textbf{No explicit required-columns list}: The written spec does not maintain a “minimum required columns” list; the implementation throws hard errors when referenced fields are missing/invalid.

\item \textbf{PD column for replenishment tests}: Tape contains a loan-level column PD interpreted as 1-year PD in [0,1], used for WAPD constraints.

\item \textbf{PD validation (implicit by global rule)}: Missing/invalid PD values -> hard error.

\item \textbf{Default event definition (v1)}: Default is modelled as a single generic hazard-driven event; no bankruptcy/failure-to-pay subtypes.

\item \textbf{Protection Start is user input}: User inputs Protection Start Date.

\item \textbf{Protection Start adjustment}: Apply Modified Following + selected calendar(s) to Protection Start.

\item \textbf{Coverage window uses default date}: Covered iff $\text{ProtStart} \le \tau \le \text{ProtEnd}$ (inclusive; default date based).

\item \textbf{Notice after ProtEnd allowed}: Covered default $\tau$ may have NoticeDate > ProtEnd; write-down still recognized (capped by Legal Final).

\item \textbf{NoticeDate default-lag stays fixed}: $\text{NoticeDateRaw} = \tau + 1 \text{ calendar month}$ (fixed, always).

\item \textbf{Notice date adjustment method}: Use Preceding with the selected joint calendar(s).

\item \textbf{Adjust then cap (restated)}: Adjust Notice and Legal Final first, then cap Notice $\le$ Legal Final.

\item \textbf{Write-down effective timing}: Write-down effective at start-of-day Notice Date; premium excludes Notice Date.

\item \textbf{Premium stop on notice (restated)}: Premium continues until Notice Date (exclusive), not until default date.

\item \textbf{Premium base unchanged until notice}: Premium accrues on the full pre-write-down notional until Notice Date.

\item \textbf{Copula simulation unit}: One latent variable and one default time per obligor (Debtor\_ID).

\item \textbf{One-factor draw per path}: For each MC path, draw one systematic factor $Y \sim N(0,1)$.

\item \textbf{Idiosyncratic draws}: For each obligor $i$ in a path, draw $\varepsilon_i \sim N(0,1)$ independent across obligors conditional on $Y$.

\item \textbf{Latent variable formula}: 
\[ X_i = \sqrt{\rho} \cdot Y + \sqrt{1-\rho} \cdot \varepsilon_i \]

\item \textbf{rho input}: User inputs a single constant $\rho$ applied to all obligors.

\item \textbf{rho validation}: Hard validate $0 \le \rho < 1$ otherwise hard error.

\vspace{1em}
\noindent Locked parameters 126–155
\vspace{0.5em}

\item \textbf{Uniform mapping from latent variable}: 
\[ U_i = \Phi(X_i) \]

\item \textbf{Default-time condition for inversion}: Compute $\tau_i$ by solving $S(\tau_i) = U_i$ (not $1-S(\tau)=U$).

\item \textbf{Default-time inversion method}: Analytic inversion using your piecewise-constant forward hazards $\lambda_k$ by tenor interval (no numerical root finding).

\item \textbf{Path computation style}: Vectorized per path—draw $Y$, draw all $\varepsilon_i$, compute all $X_i, U_i, \tau_i$.

\item \textbf{Monte Carlo RNG type}: Pseudo-random RNG (standard random numbers).

\item \textbf{Random seed convention}: Always use a fixed default seed (no user seed input).

\item \textbf{Monte Carlo path count}: Fixed default number of paths = 50,000 (no user control).

\item \textbf{Variance reduction}: None (no antithetic variates, no Sobol).

\item \textbf{No coverage after Protection End}: Defaults with $\tau > \text{Protection End}$ are not covered (restated as protection cutoff rule).

\item \textbf{Protection End boundary inclusive}: If $\tau = \text{Protection End}$, it is covered (restated).

\item \textbf{Replenishment exists}: Deal has replenishment = Yes.

\item \textbf{Replenishment end date input}: User inputs Replenishment End Date.

\item \textbf{Replenishment availability test}: If Replenishment End Date $\le$ as-of date -> treat as no replenishment remaining (skip replenishment logic).

\item \textbf{Replenishment modelling approach}: Synthetic replenishment up to cap while enforcing constraints (no candidate-loan selection).

\item \textbf{Replenishment cap definition}: 
\[ \text{Cap} = N_{\text{ref}} = \sum \text{Outstanding\_Principal\_Amount} \]
(deal-level cap; not our tranche notional).

\item \textbf{Replenishment timing}: Replace exposure lost due to amortization/maturity; do not replace exposure lost due to defaults.

\item \textbf{Replenishment granularity}: Apply replenishment on the event-date grid (not daily), topping up for amort/maturity reductions since the prior event date.

\item \textbf{Defaulted claims not replenished}: Defaulted reference claims are not replaced and are excluded from pool metrics starting at Notice Date.

\item \textbf{Synthetic replenishment characteristics}: Replenished exposure is assumed to match the current non-defaulted pool characteristics at that time (neutral top-up).

\item \textbf{Partial replenishment allowed}: If constraints would be violated by topping up fully, replenish partially up to the maximum allowed.

\item \textbf{Binding constraint rule}: Compute the maximum permissible top-up under each constraint and take the minimum (tightest constraint binds).

\item \textbf{WAPD PD source}: Use loan-level column PD (1-year PD) for WAPD calculations.

\item \textbf{WAPD formula (using current projected balances and non-defaulted pool)}: 
\[ WAPD(t) = \frac{\sum_{i \in \text{non-defaulted}} B_i(t) \cdot PD_i}{\sum_{i \in \text{non-defaulted}} B_i(t)} \]

\item \textbf{WAPD weights}: Use current projected balance $B_i(t)$ at the replenishment event date (not as-of weights).

\item \textbf{Exclude defaulted claims timing (for metrics)}: A loan is treated as defaulted for pool metrics starting Notice Date (start-of-day).

\item \textbf{WAPD no-deterioration baseline}: Compare post-replenishment WAPD to the same-date pre-replenishment pool WAPD (after amort/maturity and excluding defaulted claims).

\item \textbf{Guidelines table exists}: General portfolio guidelines are represented as a configurable table of constraints (user can change numeric thresholds; “country” refers to issuer bank’s country conceptually).

\item \textbf{Guidelines enforcement}: All guideline rows are hard constraints (breach blocks/limits replenishment via the binding constraint rule).

\item \textbf{Guidelines application under synthetic replenishment}: Evaluate constraints on the post-top-up pool where replenished exposure inherits the same characteristics as the pre-top-up non-defaulted pool.

\item \textbf{Portfolio guideline parameters are user-editable}: The model supports user-edited numeric thresholds for guideline rows (e.g., country limits, sector limits, maturity/WAL limits, PD limits, turnover/rating floors), applied consistently as hard constraints.

\vspace{1em}
\noindent \textbf{SRT Valuation Model — Locked Parameters (Sections 6–8)} \\
Compiled on 2026-02-28 (incremental decisions locked in this chat).
\vspace{0.5em}

\item \textbf{Valuation target — full junior tranche ownership}: We value 100\% of the most junior tranche (full position, not a participation). No scaling by an `our_percentage` applies in Sections 6–8.

\item \textbf{Loss allocation — junior-first until exhausted (equity behavior)}: All portfolio credit losses are allocated to the junior tranche first until the tranche is fully written down (tranche is `dead` at 0 outstanding). This is implemented as an equity/first-loss bucket (attachment effectively $A = 0$).

\item \textbf{Tranche outstanding notional definition (cap-only bucket)}: Let $N_{\text{sched}}(t)$ be the scheduled tranche notional (before defaults) and $CumLoss(t)$ cumulative portfolio credit loss amount (booked at Notice Dates). Then:
\[ N_{\text{tr}}(t) = \max(N_{\text{sched}}(t) - CumLoss(t), 0). \]
Incremental tranche loss at an event time $t$ is:
\[ \Delta TrLoss(t) = \min(\Delta Loss(t), N_{\text{tr}}(t-)). \]

\item \textbf{Scheduled tranche notional tracks scheduled/performing pool balance}: Define $PoolBal_{\text{sched}}(t)$ as the scheduled/performing pool balance from the EAD engine (contractual amortization/maturity plus replenishment top-ups, and plus/minus prepayments if enabled), but not reduced mechanically by defaults.

Define $\alpha = N_{\text{sched}}(0) / PoolBal_{\text{sched}}(0)$. Then:
\[ N_{\text{sched}}(t) = \alpha \times PoolBal_{\text{sched}}(t). \]
Defaults affect tranche only through $CumLoss(t)$, not through $PoolBal_{\text{sched}}(t)$.

\item \textbf{Cumulative loss definition (principal loss only) and recognition time}: For each defaulted debtor, principal loss amount is:
\[ Loss_i = EAD_i(\tau_i) \times LGD_i, \]
where $\tau_i$ is the debtor default time and EAD is frozen at default date.

Loss is recognized into CumLoss at Notice Date (Default Date + 1 calendar month, with business-day adjustment), with write-down effective at Notice Date start-of-day.

\item \textbf{Premium accrual base and intra-quarter segmentation}: Premium accrues on the tranche outstanding notional, with exact piecewise integration within each quarter.

For an accrual period $[T_k, T_{k+1}]$, split at any Notice Dates $t_N$ that fall inside the period. Accrue premium on each sub-interval $[u_m, u_{m+1}]$ using the prevailing outstanding notional $N_{\text{tr}}(u_m)$, where notional steps down at Notice Date start-of-day.

Day count for premium accrual is ACT/360.

\item \textbf{Premium payment timing}: Premium is paid in arrears at the period end date (typically quarter-end). Cashflow date for premium in $[T_k, T_{k+1}]$ is $T_{k+1}$.

\item \textbf{Write-down cashflow representation for funded SPV structure}: No external protection payment is assumed at default. Instead, each tranche write-down is represented as a negative note/NAV cashflow at Notice Date start-of-day:
\[ CF_{\text{wd}}(t_N) = -\Delta TrLoss(t_N). \]

\item \textbf{Redemption cashflow at Legal Final Maturity}: At Legal Final Maturity (LFM), repay remaining outstanding principal:
\[ CF_{\text{red}}(T_{\text{LFM}}) = +N_{\text{tr}}(T_{\text{LFM}}-). \]

\item \textbf{Fees and expenses}: Fees/expenses (SPV admin, trustee, etc.) are ignored in v1.

\item \textbf{Deal-level call/unwind}: Deal-level call/unwind features are ignored in v1 (deal runs to Legal Final Maturity).

\item \textbf{User toggle: stochastic loan prepayment (CPR) and replenishment interaction}: Model supports a user toggle $\text{enable\_prepayment} \in \{0,1\}$. If enabled, user inputs annual CPR ($\text{CPR\_annual} \in [0,1]$). Convert to quarterly prepayment probability:
\[ p_q = 1 - (1 - \text{CPR\_annual})^{(1/4)}. \]
Each quarter, for each live loan, draw $U \sim \text{Uniform}(0,1)$. If $U < p_q$, the loan prepays (remaining outstanding principal becomes 0 from that quarter onward).

Within the replenishment period, prepayment-driven reductions are eligible for top-off (treated like amortization/maturity). Defaults are not eligible for top-off.

\item \textbf{Discounting curve usage}: Use a single continuous discount curve $DF_{\text{ccy}}(\text{date})$ per deal currency for all cashflows: premium, write-down/NAV, and redemption. Discount factors are queried by date (curve’s own internal basis).

\item \textbf{PV computation: pathwise then average}: For each Monte Carlo path $\omega$, compute $PV_{\text{prem}}(\omega)$, $PV_{\text{wd}}(\omega)$, $PV_{\text{red}}(\omega)$ and $NPV(\omega) = PV_{\text{prem}}(\omega) + PV_{\text{wd}}(\omega) + PV_{\text{red}}(\omega)$ (note $PV_{\text{wd}}$ is negative). Then average across paths:
\[ \mathbb{E}[NPV] \approx (1/N) \sum_\omega NPV(\omega). \]

\item \textbf{Output pack (pricing) — required breakdown}: Report the following scenario-averaged outputs:
$PV_{\text{premium}}$, $PV_{\text{write\_down}}$, $PV_{\text{redemption}}$, $NPV_{\text{MTM}}$.

Also report CleanPrice and DirtyPrice per 100 of $N_{\text{tr}}(\text{as-of})$, and AccruedPremium(as-of).

\item \textbf{Dead tranche reporting convention}: If $N_{\text{tr}}(\text{as-of}) = 0$, set PV legs from as-of onward to 0, and report $NPV_{\text{MTM}} = 0$, $\text{AccruedPremium} = 0$, $\text{CleanPrice} = 0$, $\text{DirtyPrice} = 0$ (avoid NaNs).

\item \textbf{Accrued premium cutoff at as-of}: As-of is treated as start-of-day. Accrued premium is computed from last premium payment date $t_{\text{prev}}$ up to as-of $t_0$, excluding the as-of day:
\[ \text{Accrued} = s \times \sum_m \text{YearFrac}(u_m, u_{m+1}) \times N_{\text{tr}}(u_m), \]
with breakpoints $\{t_{\text{prev}}, \text{any Notice Dates in } (t_{\text{prev}}, t_0), t_0\}$.

\item \textbf{Par spread solver included (definition)}: Model includes a par spread solver. Par spread $s^*$ is defined as the spread that makes MTM NPV zero:
\[ NPV_{\text{MTM}}(s^*) = 0. \]

\item \textbf{Par spread method (closed-form PV01)}: Because NPV is linear in spread, solve par spread in closed form using PV01:
\[ NPV(s) = s \times PV01 - PV_{\text{wd}} + PV_{\text{red}}. \]
Thus:
\[ s^* = (PV_{\text{wd}} - PV_{\text{red}}) / PV01. \]
Use consistent units (e.g., if PV01 is per 1bp, $s^*$ is in bps).

\item \textbf{Validation pack (v1)}: Implement standard validation outputs:
(i) bounds: tranche loss $\le$ scheduled tranche notional; $N_{\text{tr}}(t) \ge 0$; premium stops after tranche is dead.
(ii) monotonicity: NPV increases with spread (test several spreads).
(iii) Monte Carlo convergence: compare par spread across multiple seeds and/or path counts.
(iv) loss distribution: Expected Loss, VaR(99), ES(99) for tranche loss.
(v) reconciliation: expected premium / write-down / redemption by date table.

\item \textbf{Sensitivity analysis}: Sensitivity grids are deferred (not required in v1).
\end{enumerate}

\end{document}